\documentclass[10pt,a4paper]{amsart}
\usepackage[margin=19mm]{geometry}
\usepackage{amsmath,amssymb,mathtools}
\usepackage[scr=rsfs,cal=euler]{mathalfa}
\usepackage{xfrac}
\usepackage[unicode,hidelinks,bookmarks=false]{hyperref}
\newcommand{\ie}{i.e.\ }
\newcommand{\cf}{cf.\ }
\newcommand{\resp}{respectively\ }
\newcommand{\Subsection}{\subsection}
\providecommand{\nobreakdash}{\nobreak\hbox{-}}
\setlength{\emergencystretch}{2em}
\multlinegap0pt
\usepackage{bm}
\usepackage{bbm}
\usepackage{dsfont}
\usepackage{xspace}
\usepackage{cancel}
\usepackage{mathtools}
\usepackage{amsthm}
\makeatletter
\@ifundefined{theorem}{\newtheorem{theorem}{Theorem}}{}
\@ifundefined{lemma}{\newtheorem{lemma}[theorem]{Lemma}}{}
\@ifundefined{proposition}{\newtheorem{proposition}[theorem]{Proposition}}{}
\makeatother
\title[Demotions of ideals in commutative rings with applications t]{Demotions of ideals in commutative rings with applications to normally torsion-freeness}
\author{Mehrdad Nasernejad, Jonathan Toledo}
\date{Source: arXiv:2510.16538v1 (CC BY 4.0)}
\begin{document}
\maketitle

\noindent\emph{Source and licence.} Demotions of ideals in commutative rings with applications to normally torsion-freeness. Version arXiv:2510.16538v1 (CC BY 4.0), distributed under the Creative Commons Attribution 4.0 International licence, as recorded in the arXiv licence metadata for this exact version and shown on the abstract page https://arxiv.org/abs/2510.16538v1. Full rights notice: licenses/arxiv-2510.16538-CC-BY-4.0.txt.\par\smallskip

\noindent\textit{Editorial scope.} Counted unit: the Proposition whose source label is \texttt{Summation} in the article (source0.tex lines 560--564, source label \texttt{Summation}), delivered together with the article's own complete original proof of it (source0.tex lines 565--624, one \texttt{proof} environment of 60 lines). No mathematics is written, repaired or re-derived: statement and proof are the article's own text line for line. Required context retained uncounted: LEM.SUM1 (lines 515--525). Every definition, display and label of those lines is retained unchanged. Omitted from this package and not redistributed: 1--514, 526--559, 625--1370. Nothing inside a retained excerpt is omitted or altered: every retained body is delivered exactly as the article's lines are written, so no omission receipt of any kind accompanies this package. Citations: the retained excerpts carry no citation command at all, so no bibliography entry of the article is transcribed and no reference is redistributed. Reference adaptation: none was needed and none was made; every reference inside the delivered unit resolves inside it, so no unresolved reference or citation is produced. Printed slips kept literal: no printed word, symbol, hypothesis or number of the counted statement or of its proof was altered. Layout and class: the article's own class is \texttt{amsart} with packages including amsmath, amsthm, amscd, amsfonts, amssymb, enumerate, verbatim, newlfont, calc, graphicx, color, hyperref, tikz, doi; the wrapper is the lane's pinned \texttt{amsart} editorial wrapper at 10pt on A4, which restates the theorem-like environments the retained text needs and adds the article's own macro definitions that the retained text needs (none), with extra wrapper packages (bm, bbm, dsfont, xspace, cancel, mathtools). The delivered numbering is the unit's own. Assets: no retained span contains a figure environment, a graphics-inclusion command, a PDF-page inclusion, a table or a TikZ picture; no image file is copied into this package, the package declares no asset dependency and no image file is redistributed with it. Licence: version arXiv:2510.16538v1 of this article is distributed under the Creative Commons Attribution 4.0 International licence, as recorded in the arXiv licence metadata for this exact version and shown on the abstract page https://arxiv.org/abs/2510.16538v1. A full rights notice is delivered at licenses/arxiv-2510.16538-CC-BY-4.0.txt. Characters: every retained excerpt is pure ASCII, so the delivered engine, XeLaTeX with its default Latin Modern text font, reports no missing character.
\begin{lemma} \label{LEM.SUM1}
Let $J_1 \subseteq I_1$ (respectively,  $J_2 \subseteq I_2$)  be monomial ideals  in $R=K[x_1, \ldots, x_n]$ such that 
$\mathcal{G}(I_1) \subset R_1=K[x_1, \ldots, x_m]$ and $\mathcal{G}(I_2) \subset R_2=K[x_{m+1}, \ldots, x_n]$ for some $m \geq 1$. 
Let $r,s,c,d\geq 0$ with $c+d=s$. Then 
\[
\Big(\sum_{\alpha+\beta=r+s} I_1^\alpha I_2^\beta\Big)\cap (J_1^cJ_2^d)
=
\sum_{\substack{a+b=r+s\\ a\ge c,\ b\ge d}}
\big(I_1^{a}\cap J_1^c\big)\,\big(I_2^{b}\cap J_2^d\big).
\]
\end{lemma}

\begin{proposition} \label{Summation}
 Let $J_1 \subseteq I_1$ (respectively,  $J_2 \subseteq I_2$)  be monomial ideals  in $R=K[x_1, \ldots, x_n]$ such that 
$\mathcal{G}(I_1) \subset R_1=K[x_1, \ldots, x_m]$ and $\mathcal{G}(I_2) \subset R_2=K[x_{m+1}, \ldots, x_n]$ for some $m \geq 1$. 
 If  $J_1$ is  a demotion of $I_1$ (respectively,  $J_2$ is  a demotion of $I_2$), then   $J_1+J_2$ is  a demotion of $I_1+I_2$. 
\end{proposition}

\begin{proof}
Assume that   $J_1$ is  a demotion of $I_1$ (respectively,  $J_2$ is  a demotion of $I_2$). For simplicity of notation, let us define $J := J_1 + J_2$ 
and $I := I_1 + I_2$.  Fix $r,s\geq 0$.  Since the variables in \(R_1\) and \(R_2\) are disjoint, the polynomial
 ring decomposes as a tensor product $R = R_1 \otimes_K R_2.$ It is obvious that $J \subseteq I$. 
 Moreover, based on the binomial expansion theorem,  it follows that
\[
I^r = \sum_{a+b=r} I_1^a I_2^b \quad \text{and}  \quad J^s = \sum_{c+d=s} J_1^c J_2^d,
\]
where multiplication is taken inside \(R\).  By distributivity, we have 
\[
I^r J^s = \left( \sum_{a+b=r} I_1^a I_2^b \right) \left( \sum_{c+d=s} J_1^c J_2^d \right) = \sum_{\substack{a+b=r, \; c+d=s}} I_1^a I_2^b J_1^c J_2^d.
\]

Because the variables of \(R_1\) and \(R_2\) are disjoint, these factors commute and can be grouped as $I_1^a J_1^c \cdot I_2^b J_2^d \subseteq R_1 \otimes_K R_2.$
Due to  \(J_1\) is a demotion of \(I_1\), we get $I_1^a J_1^c = I_1^{a+c} \cap J_1^c$ for all \(a, c \geq 0\). Similarly, since \(J_2\) is a demotion of \(I_2\), we have
$I_2^b J_2^d = I_2^{b+d} \cap J_2^d$ for all \(b, d \geq 0\). Substituting these back,  
\[
I^r J^s = \sum_{\substack{a+b=r, \; c+d=s}} (I_1^{a+c} \cap J_1^c)(I_2^{b+d} \cap J_2^d).
\]

By virtue of  the variables are disjoint, we can readily deduce that 
 \begin{align*}
  (I_1^{a+c} I_2^{b+d}) \cap (J_1^c J_2^d) &=   (I_1^{a+c} \cap  I_2^{b+d}) \cap (J_1^c \cap  J_2^d)\\
 &= (I_1^{a+c} \cap J_1^c) \cap ( I_2^{b+d} \cap  J_2^d)\\
 &= (I_1^{a+c} \cap J_1^c) ( I_2^{b+d} \cap  J_2^d).
 \end{align*}
Thus, we can conclude the following equality 
\[
I^r J^s = \sum_{\substack{a+b=r,\; c+d=s}} (I_1^{a+c} I_2^{b+d}) \cap (J_1^c J_2^d).
\]

 It follows now from Lemma \ref{LEM.SUM1} that 
\[
\begin{aligned}
I^{r+s}\cap J^s
&=\Big(\sum_{\alpha+\beta=r+s} I_1^\alpha I_2^\beta\Big)\cap
\Big(\sum_{c+d=s} J_1^cJ_2^d\Big) \\
&= \sum_{c+d=s}
\Big[\Big(\sum_{\alpha+\beta=r+s} I_1^\alpha I_2^\beta\Big)\cap (J_1^cJ_2^d)\Big] \\
&= \sum_{\substack{c+d=s,\; a+b=r+s\\ a\ge c,\ b\ge d}}
\big(I_1^{a}\cap J_1^c\big)\,\big(I_2^{b}\cap J_2^d\big).
\end{aligned}
\]
On account of  $J_1$ is a demotion of $I_1$ and $J_2$ is a demotion of $I_2$, for $a\ge c$ and $b\ge d$,  we have
$I_1^{a}\cap J_1^c = I_1^{a-c}J_1^c$ and $I_2^{b}\cap J_2^d = I_2^{b-d}J_2^d$. 
Thus, we obtain
\[
I^{r+s}\cap J^s = \sum_{\substack{c+d=s,\; a+b=r+s\\ a\ge c,\ b\ge d}}
I_1^{a-c}J_1^c \; I_2^{b-d}J_2^d.
\]
Let $a':=a-c$ and $b':=b-d$. Then $a'+b'=(a+b)-(c+d)=(r+s)-s=r$, and so we get the following equalities 
\[
I^{r+s}\cap J^s = \sum_{\substack{a'+b'=r, \; c+d=s}}
I_1^{a'}I_2^{b'}  J_1^cJ_2^d
=\Big(\sum_{a'+b'=r} I_1^{a'}I_2^{b'}\Big)
\Big(\sum_{c+d=s} J_1^cJ_2^d\Big)
= I^rJ^s.
\]
We can at once  conclude that   \(J = J_1 + J_2\) is a demotion of \(I = I_1 + I_2\), as claimed. 
\end{proof}

\end{document}
